\chapter{Nombre dérivé, tangente et approximation}
\section*{Exercice 8.1 -- Revenir à la définition}
Calculer le nombre dérivé de $f(x)=x^2$ en $a=-2$ à partir du taux d'accroissement.

\section*{Exercice 8.2 -- Tangente}
Déterminer la tangente à $f(x)=1/x$ au point d'abscisse $2$ et vérifier qu'elle passe par ce point.

\section*{Exercice 8.3 -- Racine à l'origine}
Pourquoi $\sqrt x$ n'a-t-elle pas de dérivée finie en $0$, même en considérant seulement $h>0$ ?

\section*{Exercice 8.4 -- Valeur absolue}
Calculer les taux d'accroissement de $|x|$ en $0$ selon le signe de $h$.

\section*{Exercice 8.5 -- Approximation affine}
Approcher $\sqrt{4{,}04}$ par la tangente à $\sqrt x$ en $4$. Déterminer le signe de l'erreur.

\section*{Exercice 8.6 -- Vitesse et unités}
Une position en mètres est $s(t)=t^2+3t$, avec $t$ en secondes. Calculer vitesse moyenne sur $[2,2+h]$ et vitesse instantanée en $2$.
